% ATTEMPT_STATUS: unresolved
% RESEARCH_GENERATED: 2026-10-01; GPT-6 Astra Ultra; individual mathematical attempt
% TOP_PROBLEM: {"id": 155, "title": "Affine Translates of {0,1,3}", "category_id": 2, "category": {"id": 2, "name": "combinatorics", "display_name": "Combinatorics", "description": "Counting problems, graph theory, discrete structures.", "slug": "combinatorics", "order_index": 2, "created_at": "2026-07-31T15:26:25.671Z"}, "status": "open"}
% FIRST_POSTED: 2026-10-01
% Imported from: unsolvedmath_additional_500.tex; UnsolvedMath ID 155
% !TeX program = lualatex
% Compile twice: lualatex 155.tex
% Self-contained: no external bibliography, images, or \input files.
% Numeric section labels and visible numbers are original UnsolvedMath IDs.
\documentclass[11pt,a4paper]{article}
\usepackage[margin=24mm,headheight=14pt]{geometry}
\usepackage{fontspec}
\setmainfont{Latin Modern Roman}
\setsansfont{Latin Modern Sans}
\setmonofont{Latin Modern Mono}
\usepackage{amsmath,amssymb,mathtools}
\usepackage{unicode-math}
\setmathfont{Latin Modern Math}
\usepackage{microtype}
\usepackage{enumitem,xurl,needspace}
\usepackage[dvipsnames]{xcolor}
\usepackage{hyperref}
\hypersetup{unicode=true,colorlinks=true,linkcolor=MidnightBlue,urlcolor=MidnightBlue,
 pdftitle={UnsolvedMath Problem 155},pdfauthor={Source-annotated compilation},bookmarksnumbered=true}
\usepackage{bookmark,tocloft,titlesec,fancyhdr}
\setlength{\cftsecnumwidth}{4.2em}
\cftsetpnumwidth{2.5em}
\cftsetrmarg{4em}
\titleformat{\section}{\Large\bfseries\raggedright}{\thesection}{0.7em}{}
\pagestyle{fancy}\fancyhf{}
\fancyhead[L]{\small\sffamily UnsolvedMath research attempt}
\fancyhead[R]{\small Original catalogue IDs}
\fancyfoot[C]{\thepage}
\setlength{\parindent}{0pt}
\setlength{\parskip}{5pt plus 1pt}
\setlength{\emergencystretch}{3em}
\setcounter{tocdepth}{1}\setcounter{secnumdepth}{1}
\setlist[itemize]{leftmargin=3em,itemsep=4pt,topsep=4pt,parsep=0pt}
\allowdisplaybreaks
\newcommand{\N}{\mathbb{N}}
\newcommand{\Z}{\mathbb{Z}}
\newcommand{\Q}{\mathbb{Q}}
\newcommand{\R}{\mathbb{R}}
\newcommand{\C}{\mathbb{C}}
\newcommand{\F}{\mathbb{F}}
\newcommand{\A}{\mathbb{A}}
\newcommand{\PP}{\mathbb{P}}
\newcommand{\E}{\mathbb{E}}
\newcommand{\eps}{\varepsilon}
\newcommand{\dd}{\,\mathrm d}
\newcommand{\sref}[2]{\hyperlink{source:#1:#2}{[S#2]}}
\newif\ifshowcatalogue
\showcataloguetrue
\newcommand{\cataloguescope}[1]{\ifshowcatalogue
 \paragraph{Statement recorded in the supplied source.}
 #1\fi}
\begin{document}
% UnsolvedMath ID 155; source catalogue code GREEN-067

\setcounter{section}{154}

\section{Counting Affine Copies of a Three-Point Pattern}\label{155}

{\small\sffamily UnsolvedMath ID 155 · Combinatorics}

\subsection{Definitions and mathematical statement}

\paragraph{Shared notation.}Unless a section says otherwise, $\N=\{1,2,\ldots\}$,
$\Z$ is the ring of integers, $\Q,\R,\C$ are the rational, real, and complex
fields, $[N]=\{1,\ldots,\lfloor N\rfloor\}$, and $\log$ is the natural
logarithm. For real functions of a parameter tending to infinity,
$f\ll g$ and $f=O(g)$ mean $|f|\leq Cg$ eventually for some fixed $C$;
$f\gg g$ reverses this comparison; $f=o(g)$ means $f/g\to0$;
$f\sim g$ means $f/g\to1$; and $f\asymp g$ means both order bounds.
A subscript on $O$ or $\ll$ specifies allowed constant dependence.
The expression $n^{o(1)}$ in an upper bound means growth smaller than
$n^\epsilon$ for each fixed $\epsilon>0$. Local definitions govern any
exception, finite-field convention, or use of the same symbol for another
object. An asymptotic claim is not automatically an effective algorithm.



A nondegenerate affine copy of $\{0,1,3\}$ is $\{a,a+d,a+3d\}$ with $a,d\in\Z$ and $d\ne0$. Count distinct three-element subsets contained in $A$, rather than arbitrary repeated parametrizations. The extremal function maximizes this count over all $n$-element subsets of $\Z$, with no bound on the magnitudes of their elements. \sref{155}{1} \sref{155}{2}

\paragraph{Statement recorded in the supplied source.}
If $A$ is a set of $n$ integers, what is the maximum number of affine translates of the set $\{0, 1, 3\}$ that $A$ can contain? \sref{155}{1}

\subsection{Short English statement}

Among all sets of a given number of integers, which can contain the most translated and dilated copies of the pattern zero, one, three?

\subsection{Research attempt}

\paragraph{Provenance and scope.}
This individual attempt was generated by GPT-6 Astra Ultra on 1 October 2026. The method was to derive explicit partial results and test the first proposed extension adversarially. The original statement and sources below are retained. No claim of mathematical novelty or complete solution is made.

\paragraph{Formal counting convention.}
A copy is a distinct three-element set $\{a,a+d,a+3d\}$ with $d\ne0$. Negative $d$ is permitted. For a sorted triple $u<v<w$, the condition is that the two adjacent gaps are in ratio $1:2$ or $2:1$. These alternatives are mutually exclusive, so every qualifying triple has exactly one parameter pair $(a,d)$. The primary problem list suggests a leading constant one third; the following derivations establish a lower construction, a universal one-half upper constant, and an explicit failure of the simplest compression proof.

\paragraph{Exact interval counts.}
For $A=[n]$ and positive $d$, the starting point can be any $a$ from one to $n-3d$. Thus there are
\[
 T_+(n)=\sum_{d=1}^{\lfloor(n-1)/3\rfloor}(n-3d)
\]
positive-orientation copies. Reflection gives the same number with negative $d$, and no triple is counted in both orientations. Therefore
\[
 T([n])=2\sum_{d=1}^{\lfloor(n-1)/3\rfloor}(n-3d)
       =\frac13n^2+O(n).
\]
The calculation includes the integer endpoints exactly, rather than replacing an ordered-parameter count by an unordered-subset count without checking multiplicity. Affine images of the interval have identical counts, since nonzero dilation and translation preserve gap ratios.

\paragraph{A universal bound by middle rank.}
Write an arbitrary integer set as $a_1<\cdots<a_n$. Fix its middle point $a_i$ in a potential copy. In the orientation with shorter left gap, a chosen left point determines the right point uniquely by $a_i+2(a_i-a_j)$, and a chosen right point likewise determines the left point if it is integral. Hence there are at most $\min(i-1,n-i)$ such copies. The opposite orientation satisfies the same bound. Summing yields
\[
 T(A)\leq2\sum_{i=1}^n\min(i-1,n-i)
 =\begin{cases}
 2m(m-1),&n=2m,\\
 2m^2,&n=2m+1.
 \end{cases}
\]
In particular $T(A)\leq n^2/2+O(n)$. Each triple has a unique middle point, so the sum does not introduce an extra multiplicity. The inequality uses order but not a bound on the size of the integers, making it genuinely applicable to the unrestricted extremal problem.

\paragraph{A finite obstruction to monotone compression.}
It is tempting to move the sorted elements onto consecutive integers and claim that the number of copies cannot decrease. This is false. The set
\[
 A=\{0,1,2,3,4,6\}
\]
has seven copies:
\[
 \{0,1,3\},\{0,2,3\},\{0,2,6\},\{0,4,6\},
 \{1,2,4\},\{1,3,4\},\{3,4,6\}.
\]
Its rank compression $\{0,1,2,3,4,5\}$ has only six, as the interval formula gives. Each displayed triple is checked by its adjacent gaps, and the companion script enumerates all parameter pairs to verify completeness. Thus consecutive intervals need not be exact finite maximisers, even though they could still give the correct leading constant. The example refutes an approach, not the proposed asymptotic constant.

\paragraph{What the middle-rank bound is discarding.}
At a central element the upper bound allows almost every point on each side to participate in each of the two possible ratios. But the same surrounding points are constrained at other centres. The rank argument treats these demands independently. Reaching the conjectured one-third constant requires a global incompatibility statement between those local matchings. A putative proof that simply replaces $\min(i-1,n-i)$ by a smaller universal factor at every centre is unlikely to be valid: geometric strings of distances can make a particular centre support many alternating ratio-two links. The saving, if available, must be accumulated across centres rather than imposed identically at each one.

The finite compression example makes this dependence concrete. The final point six creates two long-scale triples at zero while retaining a short-scale triple through three. Moving it to five destroys some of these coincidences while creating others; local distance shortening has no monotone effect on a prescribed ratio.

\paragraph{Verification and remaining gap.}
For $n<3$ both formulas give zero where appropriate, and for $n=3$ the upper bound permits the one possible qualifying triple. Equal factors or zero dilation are excluded by the three-distinct-points convention. The exact interval lower count, the middle-rank upper bound, and the seven-copy example are proved or exhaustively checked. They leave a gap between leading constants one third and one half. The unresolved task is a uniform asymptotic upper bound with the lower construction's constant, or a family beating it by a positive quadratic proportion.

\subsection{Sources}

{\small

\begin{itemize}

\item[\hypertarget{source:155:1}{S1}] ULAM.AI, \emph{UnsolvedMath}, supplied \texttt{problems.json}, record 155 (GREEN-067), statement and background. The embedded literature-review date is 2026-08-17; that is the archive’s review date, not a fresh certification. Dataset landing page: \url{https://huggingface.co/datasets/ulamai/UnsolvedMath}.

\item[\hypertarget{source:155:2}{S2}] Ben Green, 100 Open Problems (author’s maintained problem list). \url{https://people.maths.ox.ac.uk/greenbj/papers/open-problems.pdf}. The author-hosted document was inspected on 27 September 2026; the archive’s local code is not a problem-number locator in this paper.

\item[\hypertarget{source:155:3}{S3}] UnsolvedMath, GREEN-067 Affine Translates of {0,1,3}. \url{https://www.unsolvedmath.com/problems/GREEN-067}. Bibliographic information imported from the supplied record.

\end{itemize}

}

\label{end:155}
\end{document}
